\section{问题二模型的建立与求解}

\subsection{建模思路}
问题二中，小区负载与光伏出力不再具有固定的日内规律，日前计划与实际运行之间存在预测偏差，
同时供电不足需以正常电价的 5 倍进行紧急购电。因此，在问题一优化模型的基础上，
引入负载与光伏预测、预测余量及日内实时纠偏机制，构建“日前预测—计划优化—实时纠偏—跨日滚动”的调控模型。

首先，对 2025 年历史负载与光伏数据进行时序特征分析：小区负载呈现显著的星期周期特征，
光伏出力具有较强的时间连续性；据此采用上周同星期日的实际数据预测当日各 $10\,$min 时段的负载与光伏出力。
考虑到供电不足产生的紧急购电成本远高于正常购电成本，在预测结果上引入安全余量，
并通过不同预测方法与余量参数组合的全年滚动仿真确定修正方案。

其次，以修正后的负载与光伏预测值作为每日 0:00 可获知的信息，
在问题一线性规划模型的基础上制定当天计划购电及储能充放电方案；
实际运行过程中，将日前计划与当日实际负载、光伏出力逐时段比较，利用储能优先吸收预测偏差，
储能调节不足以满足负荷需求时再通过紧急购电补足电力缺口；
每日实际运行结束后的储电量作为次日优化的初始状态，从而实现全年连续滚动调度。

\subsection{数据规律分析与日前预测模型}
为在 0:00 制定当天计划购电策略，需要在不使用当日实际运行信息的前提下，对当日负载及光伏出力进行预测。
为此，对附件 2 中 2025 年实际负载和光伏数据进行统计分析，分别计算各日总负载、
总光伏发电量及不同星期类型下的负载分布，结果如图~\ref{fig:q2-data} 所示。

\begin{figure}[!htbp]
\centering
\includegraphics[width=0.48\textwidth]{Q2_weekday_box.pdf}\hfill
\includegraphics[width=0.48\textwidth]{Q2_annual_load_pv.pdf}
\caption{2025 年负载与光伏出力时序特征：（左）按星期分组的日负载电量箱线图；
（右）全年逐日负载与光伏电量}
\label{fig:q2-data}
\end{figure}

由图~\ref{fig:q2-data}（左）可知，小区日负载具有显著的星期周期特征：周日至周四的日负载整体处于较高水平，
而周五、周六明显降低，呈现较稳定的星期型双状态结构。因此，若直接采用前一日数据预测当日负载，
在周四至周五、周六至周日等状态切换时容易产生较大的预测偏差。

由图~\ref{fig:q2-data}（右）可知，全年负载在保持星期周期波动的同时还存在一定的季节变化；
光伏日发电量则呈现较明显的季节性趋势，并具有较强的时间连续性——相邻日光伏发电量的自相关系数约为 0.964\cite{fpp3} 第 2 章，
说明近期历史光伏数据对当日出力具有较强的参考价值。

综合上述特征，为兼顾预测效果、模型可解释性及信息可得性，本文采用上周同星期日预测法（w1）\cite{fpp3} 第 5.2 节 构造日前预测。
对于第 $d$ 日第 $t$ 个 $10\,$min 时段，有
\begin{equation}
\hat{L}_{d,t}=L_{d-7,t},\qquad \hat{P}^{pv}_{d,t}=P^{pv}_{d-7,t}
\end{equation}
其中 $\hat{L}_{d,t}$ 和 $\hat{P}^{pv}_{d,t}$ 分别表示第 $d$ 日第 $t$ 时段的预测负载和预测光伏功率，
$L_{d-7,t}$ 和 $P^{pv}_{d-7,t}$ 分别表示上周同星期日对应时段的实际负载和光伏功率。
该预测方法仅使用制定计划时已经获得的历史数据，不引入当日实际运行信息，从而避免未来信息泄漏；
同时能够保留负载的星期周期结构以及光伏出力的近期变化特征，为后续购电优化提供预测输入。
为进一步检验预测方法的适用性，本文同时构建持续性预测、近 2 个和近 4 个同星期日均值预测以及 7 日移动平均预测等候选方案，
并结合全年运行总费用进行比较，具体参数选择结果见下节。

\subsection{预测余量设置与参数寻优}
由于日前预测不可避免地存在误差，若直接根据预测值制定购电计划，
当实际净负荷高于预测值时可能出现供电缺口，并产生正常电价 5 倍的紧急购电费用；
相比之下，适度提高日前计划购电量所增加的正常购电成本较低。
因此，预测偏低与预测偏高所造成的经济损失具有明显的\textbf{不对称性}，需要在日前预测结果中预留一定的安全余量。
为此，引入负载余量系数 $\gamma$ 和光伏余量系数 $\delta$，对日前预测结果进行修正：
\begin{equation}
L^{*}_{d,t}=\gamma\hat{L}_{d,t},\qquad P^{pv,*}_{d,t}=\delta\hat{P}^{pv}_{d,t}
\end{equation}
其中 $\gamma\ge 1$、$0<\delta\le 1$。通过适当提高预测负载并降低预测光伏出力，
使日前优化采用较为保守的净负荷预测，从而降低实际运行过程中因供电不足而产生高额紧急购电费用的风险。

为确定日前预测方法及安全余量参数，本文分别选取持续性预测、上周同星期日预测（w1）、
近 2 个及近 4 个同星期日均值预测（w2、w4）和 7 日移动平均预测（ma7）作为候选预测方法，
在此基础上引入负载余量系数 $\gamma$ 和光伏余量系数 $\delta$，采用网格搜索法设置多组候选参数组合；
对于每一种预测方法及参数组合，均按照“日前计划—实际运行—储能纠偏—跨日状态传递”的过程进行全年滚动仿真，
并计算相应的计划购电费用与紧急购电费用，以二者之和作为方案评价指标。
比较结果表明，上周同星期日预测（w1）能够较好地刻画负载的星期周期特征，
其全年运行费用明显低于持续性预测等基准方法；在 w1 预测基础上进一步设置安全余量后，
紧急购电造成的高额成本得到进一步控制。综合各候选参数组合的全年总费用，最终选取
\begin{equation}
\gamma=1.05,\qquad \delta=0.98
\end{equation}
作为预测余量参数，即将日前负载预测值上调 5\%、光伏预测值下调 2\%。
该参数组合在适当增加计划购电量的同时，有效降低了因预测偏差造成的 5 倍电价紧急购电风险，
实现了计划购电成本与供电不足风险之间的平衡。

\subsection{日前计划购电优化模型}
在获得经安全余量修正后的负载与光伏预测结果后，以其作为每日 0:00 制定购电计划时的已知信息，
在问题一线性规划模型的基础上，建立第 $d$ 日的日前计划购电优化模型。

对于第 $d$ 日第 $t$ 个时段，微网的计划功率平衡关系为
\begin{equation}
p_{d,t}+P^{pv,*}_{d,t}+u^{dis}_{d,t}=L^{*}_{d,t}+u^{ch}_{d,t}+q_{d,t},\qquad t=1,2,\dots,144
\end{equation}
其中 $L^{*}_{d,t}$ 与 $P^{pv,*}_{d,t}$ 分别为经余量修正后的负载预测功率和光伏预测功率。
储能设备的状态转移关系为
\begin{equation}
S_{d,t}=S_{d,t-1}+\eta_c u^{ch}_{d,t}\Delta t-\frac{u^{dis}_{d,t}\Delta t}{\eta_d}
\end{equation}
并满足储电量及充放电功率约束
\begin{equation}
1200\le S_{d,t}\le 10800,\qquad 0\le u^{ch}_{d,t}\le 5000,\qquad 0\le u^{dis}_{d,t}\le 5000 .
\end{equation}

需要说明的是，问题二同时涉及\textbf{日前计划状态}与\textbf{实际运行状态}。
每日 0:00 制定购电计划时，以该时刻储能设备的实际储电量 $S_{d,0}$ 作为当日优化的初始状态，
并根据修正后的负载与光伏预测值规划全天储能运行轨迹。
为避免模型通过消耗已有储能电量人为降低当日计划购电费用，在日前计划层面设置首末储电量相等约束：
\begin{equation}
S^{plan}_{d,144}=S_{d,0}
\end{equation}
该约束仅作用于 0:00 制定的计划储能轨迹，并不要求实际运行结束后的储电量必须回到日初水平。
实际运行过程中，由于真实负载及光伏出力与日前预测存在偏差，储能设备将参与实时纠偏，
因此实际日末储电量 $S^{real}_{d,144}$ 可能与计划值存在差异。该实际日末状态继续作为下一日优化的初始状态，即
\begin{equation}
S_{d+1,0}=S^{real}_{d,144}
\end{equation}
由此实现储能状态在相邻日期之间的连续传递。

此外，由于每日 0:00 制定计划时无法获得当日真实负载及光伏出力，
因此无法预先确定实际运行过程中产生的紧急购电量。在给定预测余量参数的条件下，
日前优化以当日计划购电费用最小为目标：
\begin{equation}
\min\ C^{plan}_{d}=\sum_{t=1}^{144} c_t\,p_{d,t}\,\Delta t
\end{equation}
由此得到第 $d$ 日 144 个 $10\,$min 时段的计划购电功率及储能充放电计划。
实际运行过程中产生的预测偏差及紧急购电费用则在日内纠偏阶段根据真实负载和光伏出力计算；
预测余量参数的优选以全年计划购电费用与紧急购电费用之和作为评价指标。
模型中的储能容量、充放电效率及最大充放电功率等运行参数与问题一保持一致。

实际运行过程中产生的紧急购电费用，根据当日真实负载、光伏出力以及储能实时纠偏结果计算。
定义第 $d$ 日紧急购电费用为
\begin{equation}
C^{em}_{d}=\sum_{t=1}^{144} 5\,c_t\,e_{d,t}\,\Delta t
\end{equation}
其中 $e_{d,t}$ 为第 $d$ 日第 $t$ 时段的紧急购电功率。由此，全年实际购电总费用为
\begin{equation}
C_{total}=\sum_{d}\left(C^{plan}_{d}+C^{em}_{d}\right).
\end{equation}
需要指出的是，$C_{total}$ 并非每日 0:00 日前线性规划中直接优化的目标，
而是用于评价不同预测方法及余量参数下全年调控策略的实际经济效果——
通过对不同预测方法与余量参数组合进行全年滚动仿真，以 $C_{total}$ 作为方案评价指标，
从而实现计划购电成本与预测偏差所引致紧急购电风险之间的权衡。

\subsection{日内实时纠偏策略}
日前计划制定后，由于实际负载与光伏出力可能偏离预测值，实际运行过程中的电力供需状态也将与日前计划产生偏差。
为直观刻画预测误差的方向，首先定义第 $d$ 日第 $t$ 时段的净负荷预测偏差为
\begin{equation}
\Delta P_{d,t}=\left(L^{real}_{d,t}-P^{pv,real}_{d,t}\right)-\left(L^{*}_{d,t}-P^{pv,*}_{d,t}\right)
\end{equation}
其中第一项表示实际净负荷，第二项表示日前优化所采用的修正预测净负荷。
$\Delta P_{d,t}>0$ 表示实际净负荷高于日前预测，系统面临供电不足风险；
$\Delta P_{d,t}<0$ 则表示实际净负荷低于日前预测，系统可能出现电力富余。

然而，预测偏差并不等同于实际需要补偿的功率缺口——由于日前优化已经确定了计划购电及储能充放电方案，
实际运行时还需结合这些计划决策判断系统的真实供需状态。因此，进一步定义执行日前计划后的实时功率缺口：
\begin{equation}
D_{d,t}=L^{real}_{d,t}-\left[p^{plan}_{d,t}+P^{pv,real}_{d,t}+u^{dis,plan}_{d,t}-u^{ch,plan}_{d,t}\right]
\end{equation}
其中 $D_{d,t}$ 表示在保持日前购电计划和原储能计划不变的情况下，第 $t$ 时段实际负荷与可供电力功率之间的差额。
当 $D_{d,t}>0$ 时，说明按照日前计划执行后仍存在电力缺口；当 $D_{d,t}<0$ 时，则说明系统存在电力富余。

\textbf{（1）供电不足时的纠偏策略。} 当 $D_{d,t}>0$ 时，考虑到全天计划购电采用照付不议方式结算，
本文保持日前计划购电量不变，实时纠偏严格按照以下优先级顺序进行：
\begin{center}
\fbox{减少计划充电 $\rightarrow$ 增加储能放电 $\rightarrow$ 紧急购电}
\end{center}
首先减少该时段原计划用于储能充电的功率，将相应电能直接用于满足负荷；
若仍存在缺口，则在最大放电功率和最低储电量约束允许的范围内增加储能放电；
若储能设备达到调节极限后仍无法满足负荷需求，则对剩余缺口进行紧急购电，并按照对应时段正常电价的 5 倍计费。

\textbf{（2）电力富余时的纠偏策略。} 当 $D_{d,t}<0$ 时，调节同样按照下述优先级顺序进行：
\begin{center}
\fbox{减少计划放电 $\rightarrow$ 增加储能充电 $\rightarrow$ 弃置剩余电量}
\end{center}
首先减少该时段原计划的储能放电；若仍存在富余电力，则在最大充电功率及储电量上限允许的范围内增加储能充电；
若储能设备达到调节极限后仍存在剩余电力，则将其作为弃置电量处理。

上述两类纠偏策略的完整判断过程与执行顺序如图~\ref{fig:q2-flow} 所示：
系统先根据功率缺口 $D_{d,t}$ 的正负判断电力供需状态，再依次尝试成本由低到高的调节手段，
只有在储能调节能力耗尽后，才对剩余缺口动用紧急购电或对剩余富余电力作弃光处理。

\begin{figure}[!htbp]
\centering
\footnotesize
\begin{tikzpicture}[
  >=Stealth,
  node distance=7mm,
  io/.style={rectangle, rounded corners=2mm, draw=black!75, fill=gray!12, align=center,
             text width=66mm, minimum height=9mm, inner sep=2pt},
  proc/.style={rectangle, draw=black!70, fill=blue!6, align=center,
             text width=66mm, minimum height=8mm, inner sep=2pt},
  dec/.style={diamond, draw=black!75, fill=yellow!20, aspect=3.2, align=center, inner sep=1pt},
  act/.style={rectangle, rounded corners=1mm, draw=black!70, fill=green!10, align=center,
             text width=50mm, minimum height=8mm, inner sep=2pt},
  line/.style={draw, -Latex, thick}
]
  \node[io]   (in) {第 $t$ 个 $10\,$min 时段：按日前计划执行后计算功率缺口 $D_{d,t}$};
  \node[dec]  (d1) [below=of in] {$D_{d,t}>0$？};
  \node[act]  (a1) [right=12mm of d1] {（1）减少计划充电\\ （2）增加储能放电\\ （3）紧急购电（单价 $5c_t$）};
  \node[dec]  (d2) [below=of d1] {$D_{d,t}<0$？};
  \node[act]  (a2) [right=12mm of d2] {（1）减少计划放电\\ （2）增加储能充电\\ （3）弃置剩余电量};
  \node[proc] (ex) [below=of d2] {按日前计划执行};
  \node[proc] (up) [below=of ex] {更新实际储电量 $S^{real}_{d,t}$，进入下一时段};
  \draw[line] (in.south) -- (d1.north);
  \draw[line] (d1.south) -- node[left] {否} (d2.north);
  \draw[line] (d2.south) -- node[left] {否} (ex.north);
  \draw[line] (ex.south) -- (up.north);
  \draw[line] (d1.east) -- node[above] {是} (a1.west);
  \draw[line] (d2.east) -- node[above] {是} (a2.west);
  \coordinate (r1) at ([xshift=8mm]a1.east);
  \coordinate (r2) at ([xshift=8mm]a2.east);
  \draw[line] (a1.east) -- (r1);
  \draw[line] (a2.east) -- (r2);
  \draw[draw, thick] (r1) -- (r2);
  \draw[line] (r2) |- (up.east);
\end{tikzpicture}
\caption{日内实时纠偏决策流程（第 $d$ 日第 $t$ 时段）}
\label{fig:q2-flow}
\end{figure}

\textbf{（3）实际储能状态更新。} 完成第 $t$ 时段的实时纠偏后，根据实际充放电功率更新储能状态：
\begin{equation}
S^{real}_{d,t}=S^{real}_{d,t-1}+\eta_c u^{ch,real}_{d,t}\Delta t-\frac{u^{dis,real}_{d,t}\Delta t}{\eta_d}
\end{equation}
因此，实际储能轨迹会随着负载及光伏预测误差不断偏离日前计划轨迹。
每日运行结束后，将实际日末储电量作为次日 0:00 的初始状态：
\begin{equation}
S_{d+1,0}=S^{real}_{d,144}
\end{equation}
由此实现全年储能状态的连续滚动。

\subsection{全年滚动求解流程}
问题二需要连续制定 2025 年全年的日前购电策略，因此各日优化并非相互独立，
而是通过储能设备的实际日末状态实现跨日衔接。根据题目给定的初始条件，
设置 2025 年 1 月 1 日 0:00 储能设备的初始储电量为
\begin{equation}
S^{real}_{1,0}=6000\ \text{kWh}.
\end{equation}
随后按照“\textbf{日前预测—余量修正—计划优化—实时纠偏—状态传递}”的方式逐日滚动求解，具体流程如下：

\begin{itemize}[leftmargin=2em, itemsep=0.2em]
\item \textbf{日前预测：}在第 $d$ 日 0:00，仅利用此前已经获得的历史数据，根据上周同星期日对应 144 个时段的实际负载与光伏出力，得到当日原始预测值 $\hat{L}_{d,t}$、$\hat{P}^{pv}_{d,t}$。
\item \textbf{预测余量修正：}根据参数寻优结果，对原始预测进行安全余量修正：$L^{*}_{d,t}=1.05\hat{L}_{d,t}$，$P^{pv,*}_{d,t}=0.98\hat{P}^{pv}_{d,t}$。
\item \textbf{日前计划优化：}以当日 0:00 的实际储电量 $S^{real}_{d,0}$ 为初始状态，将修正后的负载及光伏预测值输入日前线性规划模型，求得 144 个时段的计划购电功率 $p^{plan}_{d,t}$、计划充放电功率以及计划储能轨迹。
\item \textbf{日内实际运行与实时纠偏：}将附件 2 中当日实际负载及光伏出力逐时段代入运行过程，并与日前计划进行比较。当系统出现供电缺口时，依次通过减少计划充电、增加储能放电及紧急购电进行补偿；当系统出现电力富余时，依次通过减少计划放电、增加储能充电及弃置剩余电量进行调节。每完成一个 $10\,$min 时段，即根据实际充放电情况更新储电量。
\item \textbf{跨日状态传递：}当第 $d$ 日 144 个时段全部运行结束后，将实际日末储电量传递至下一日：$S^{real}_{d+1,0}=S^{real}_{d,144}$。
\end{itemize}

随后重复上述过程，直至完成全年逐日滚动模拟。由于上周同星期日预测需要至少 7 天历史数据，
程序从 2025 年 1 月 1 日给定的 6000 kWh 初始储电量开始运行，并利用前期数据完成储能状态的连续传递；
根据题目要求，最终统计 \textbf{2025 年 2 月 1 日至 12 月 31 日共 334 天}的计划购电量、计划购电费用及紧急购电费用。

\subsection{求解结果与分析}
基于上述预测方法、日前优化模型及日内实时纠偏策略，对 2025 年全年进行逐日滚动求解，
并按题目要求统计 2025 年 2 月 1 日至 12 月 31 日的运行结果。计算得到，报告期内计划购电总量为 22320727.5\,kWh，
计划购电费用为 13676311.49 元；受负载及光伏预测误差影响，部分时段在储能完成实时纠偏后仍存在供电缺口，
全年紧急购电总量为 145133.8\,kWh，对应紧急购电费用为 870192.48 元。
因此，报告期内微网实际购电总费用为
\begin{equation}
C_{total}=13676311.49+870192.48=14546503.97\ \text{元},
\end{equation}
即约 1454.65 万元。334 个统计日中，共有 102 天发生紧急购电。
总体而言，紧急购电量仅占外网购电量的较小部分，
说明日前预测与安全余量能够覆盖大部分负荷需求偏差，而储能系统进一步承担了日内预测误差的实时调节作用。
全年每日紧急购电量的时间分布如图~\ref{fig:q2-emerg-year} 所示。

\begin{figure}[!htbp]
\centering
\includegraphics[width=\textwidth]{Q2_emergency_year.pdf}
\caption{2025 年 2--12 月每日紧急购电量分布（虚线为题目指定的四个日期）：
全年共 102 天发生紧急购电，合计 14.5 万\,kWh，集中于 6 月与 12 月的负载状态切换期}
\label{fig:q2-emerg-year}
\end{figure}

\textbf{（1）指定日期购电结果}

根据题目要求，进一步统计 2025 年 3 月 20 日、6 月 21 日、9 月 23 日和 12 月 21 日的计划购电量、
计划购电费用及紧急购电情况，结果如表~\ref{tab:q2-plan}--\ref{tab:q2-emerg} 所示。

\begin{table}[!htbp]
\caption{四个指定日期微网在指定时间段的计划购电量及全天的购电量和购电费}\label{tab:q2-plan}
\centering
\small
\setlength{\tabcolsep}{3.5pt}
\begin{tabular}{lccccccrr}
\toprule
日期 & 10:00--10:10 & 12:00--12:10 & 14:00--14:10 & 16:00--16:10 & 18:00--18:10 & 20:00--20:10 & 全天/kWh & 购电费/元 \\
\midrule
2025-03-20 & 0.0000 & 822.2488 & 0.0000 & 493.8229 & 759.4335 & 0.0000 & 78872.98 & 48644.97 \\
2025-06-21 & 0.0000 & 0.0000 & 0.0000 & 139.1066 & 451.7126 & 0.0000 & 38441.60 & 21995.11 \\
2025-09-23 & 0.0000 & 627.7958 & 0.0000 & 401.8831 & 833.4051 & 9.9357 & 70972.59 & 44123.14 \\
2025-12-21 & 0.0000 & 1023.8313 & 0.0000 & 857.8885 & 753.0828 & 0.0000 & 100592.96 & 65063.32 \\
\bottomrule
\end{tabular}
\end{table}

\begin{table}[!htbp]
\caption{四个指定日期储能设备在指定时间段的充电量（单位：kWh）}\label{tab:q2-ch}
\centering
\small
\begin{tabular}{lcccccc}
\toprule
日期 & 0:00--4:00 & 4:00--8:00 & 8:00--12:00 & 12:00--16:00 & 16:00--20:00 & 20:00--24:00 \\
\midrule
2025-03-20 & 119.4939 & 1317.6609 & 5604.5696 & 1321.6289 & 16.6683 & 10054.6548 \\
2025-06-21 & 1044.5943 & 1662.7434 & 9290.3800 & 0.0000 & 63.5124 & 9782.9959 \\
2025-09-23 & 1299.5142 & 368.7525 & 4237.2075 & 6815.3460 & 336.7457 & 10051.6290 \\
2025-12-21 & 115.7516 & 953.0670 & 5057.4378 & 7261.3994 & 0.0000 & 10112.5625 \\
\bottomrule
\end{tabular}
\end{table}

\begin{table}[!htbp]
\caption{四个指定日期储能设备在指定时间段的放电量（单位：kWh）}\label{tab:q2-dis}
\centering
\small
\begin{tabular}{lcccccc}
\toprule
日期 & 0:00--4:00 & 4:00--8:00 & 8:00--12:00 & 12:00--16:00 & 16:00--20:00 & 20:00--24:00 \\
\midrule
2025-03-20 & 96.7901 & 3157.2466 & 3265.5568 & 254.7228 & 4501.1868 & 3656.5848 \\
2025-06-21 & 7968.5140 & 1749.6373 & 0.0000 & 0.0000 & 4916.3962 & 3059.2756 \\
2025-09-23 & 1183.1489 & 6151.4626 & 2387.9403 & 440.9339 & 6648.7788 & 1954.2971 \\
2025-12-21 & 132.1232 & 671.8818 & 7758.1487 & 2220.1094 & 4416.0676 & 3775.1080 \\
\bottomrule
\end{tabular}
\end{table}

\begin{table}[!htbp]
\caption{四个指定日期的 0:00 与 24:00 储电量（单位：kWh）及紧急购电情况}\label{tab:q2-emerg}
\centering
\small
\begin{tabular}{lccc}
\toprule
日期 & 0:00 储电量 & 24:00 储电量 & 紧急购电时段与电量 \\
\midrule
2025-03-20 & 10800.00 & 10800.00 & 无 \\
2025-06-21 & 10800.00 & 10800.00 & 无 \\
2025-09-23 & 10644.02 & 10590.56 & 无 \\
2025-12-21 & 10731.40 & 10800.00 & 8:00--8:10，5.6161\,kWh \\
\bottomrule
\end{tabular}
\end{table}

\textbf{（2）典型日实际运行结果分析}

为进一步说明日前计划与日内实时纠偏策略的实际运行过程，选取 2025 年 12 月 21 日作为典型日，
对当天电价、负载预测与实际负载、计划购电以及储能运行状态进行分析，结果如图~\ref{fig:q2-plan-actual} 和图~\ref{fig:q2-day} 所示。

\begin{figure}[!htbp]
\centering
\includegraphics[width=\textwidth]{Q2_plan_vs_actual.pdf}
\caption{2025 年 12 月 21 日日前计划与实际运行对照：(a) 分时电价与计划购电功率（竖线处标注题目指定的 6 个时段对应的购电量，单位：kWh）；
(b) 日前计划储电量与实际储电量轨迹（虚线为储电量上下限 1200/10800\,kWh）}
\label{fig:q2-plan-actual}
\end{figure}

由图~\ref{fig:q2-plan-actual}(a) 可知，当日分时电价呈明显的双峰特征，计划购电功率并未完全跟随电价的高低变化：
在 8:00--11:00 与 17:00--21:00 两个高电价时段，计划购电功率被压至较低水平且时段内购电量较小；
而在 0:00--8:00 等低电价时段则维持较高的购电功率，用于满足负荷并为储能补充电量，
从而实现电能在时间维度上的“低储高放”，降低全天购电成本。

由图~\ref{fig:q2-plan-actual}(b) 可知，日前计划储电量（绿色虚线）与实际储电量（实线）总体走势一致，
但实际轨迹在预测偏差累积的时段明显偏离计划值：10:30--11:00 与 21:00 前后实际储电量一度接近运行下限 1200\,kWh，
表明该时段储能调节能力已被充分调用；此后实际储电量在低价时段重新充至上限 10800\,kWh 附近，
全天始终未越限运行。这也解释了当日仅在 8:00--8:10 出现 5.62\,kWh 短时紧急购电的原因——
绝大部分预测偏差已由储能的实时纠偏吸收，仅剩余缺口需要以高成本紧急购电兜底。

\begin{figure}[!htbp]
\centering
\includegraphics[width=\textwidth]{Q2_representative_day.pdf}
\caption{2025 年 12 月 21 日微网实际运行结果：(a) 分时电价；(b) 实际负载、w1 预测负载与计划购电功率；
(c) 实际光伏与 w1 预测光伏；(d) 储电量轨迹（虚线标注 8:00--8:10 紧急购电 5.62\,kWh）}
\label{fig:q2-day}
\end{figure}

由图~\ref{fig:q2-day} 可知，当日实际负载总体保持明显的分时变化特征，上周同星期日预测能够较好地跟踪实际负载的整体变化趋势，
但局部时段仍存在一定预测偏差。计划购电功率并非简单跟随负载变化，而是在满足预测负荷需求的基础上结合分时电价与储能充放电能力进行优化：
在部分高电价时段减少外网购电并利用储能放电供能，在低电价时段则增加购电并为储能补充电量，从而实现电能在不同时段之间的转移。

从储电量变化可以看出，储能设备在全天运行过程中多次经历充、放电过程，并始终保持在允许的容量范围内。
当实际负载偏离日前预测时，储能系统通过调整原计划充放电功率优先吸收预测误差；
仅在约 8:00 附近出现较大的供电缺口且储能调节能力不足时发生短时紧急购电。
随后储能状态继续随实际运行情况动态调整，体现了“日前计划$+$日内纠偏”策略对预测误差的适应能力。

\textbf{（3）结果小结}

综合全年滚动结果及典型日运行情况可知，所建立的调控策略能够利用历史数据制定日前购电计划，
并通过储能设备在实际运行阶段对预测偏差进行实时修正。当储能调节能力不足时，再利用高成本紧急购电保证负荷需求，
从而形成“日前优化—储能纠偏—紧急购电兜底”的分层调控机制。
全年总购电费用为 1454.65 万元，为后续不同策略的经济性比较及模型有效性检验提供了基础。
